% Introduction section for EpiContext paper (v7 - Breakthrough Results)
\section{Introduction}

The rapid advancement of large language models (LLMs) has catalyzed a paradigm shift in artificial intelligence: from single-turn question answering to autonomous agents capable of executing complex, multi-step tasks over extended horizons~\citep{yao2023react, park2023generative, wang2024survey}. Modern LLM agents---exemplified by systems like Claude Code, Cursor, and OpenHands---routinely engage in software engineering tasks spanning hundreds of turns, each generating observations, tool calls, and reasoning traces that accumulate into massive context histories~\citep{jimenez2024swe, zhou2023webarena, liu2023agentbench}.

This accumulation creates a fundamental tension. On one hand, comprehensive context is essential for coherent long-horizon reasoning: agents must remember architectural decisions made dozens of turns ago, track evolving file states, and maintain awareness of previously attempted solutions~\citep{liu2024lost}. On the other hand, context windows, while growing rapidly (from 128K to 2M+ tokens), remain finite and expensive resources. More critically, a growing body of evidence demonstrates that LLM performance degrades as context length increases---a phenomenon known as ``lost in the middle''~\citep{liu2024lost}---and that agent-generated context rapidly transforms from useful signal into distracting noise~\citep{jetbrains2025context}.

Existing approaches to this tension fall into three broad categories. \textbf{Memory compression} methods summarize or abstract conversation history into compact representations~\citep{packer2023memgpt, zhong2024memorybank, xu2025amem}. While effective at reducing token counts, these methods are inherently lossy: critical details can be irretrievably discarded, and iterative summarization introduces semantic drift~\citep{chen2025promem}. \textbf{Selective retrieval} approaches use similarity search to fetch relevant memories on demand~\citep{lewis2020rag, asai2023selfrag}, but retrieval quality depends heavily on query formulation and embedding alignment. \textbf{Dynamic tool selection} methods filter tool schemas to reduce context bloat~\citep{zhang2025autotool, anthropic2025context}, yet treat tool management independently from memory management.

We observe that all existing approaches share a common limitation: they treat context management as a \textit{static optimization problem}---compress, retrieve, or filter once, then proceed. In contrast, biological systems employ a fundamentally different strategy. In epigenetics, an organism's complete DNA sequence (genome) remains fixed, but gene expression is dynamically regulated through mechanisms like DNA methylation (gene silencing) and histone acetylation (gene activation) in response to environmental conditions~\citep{allis2015epigenetics}. This enables a single genome to produce diverse cellular phenotypes optimized for specific contexts.

Inspired by this biological principle, we introduce \EpiContext, a framework that reformulates agent context management as a \textit{dynamic epigenetic expression process}. The key insight is that an agent's complete interaction history should be treated as an immutable ``genome''---a structured context graph preserving all information---while each request payload sent to the LLM represents a context-dependent ``epigenetic expression'' regulated by three operators:

\begin{enumerate}
    \item \textbf{Memory Methylation} ($\mathcal{M}$): Silences noisy or resolved interaction blocks by reducing their activation weights, replacing them with compact summary nodes that preserve essential semantic content.
    \item \textbf{Tool Acetylation} ($\mathcal{A}$): Dynamically activates only the subset of tools and historical context relevant to the current task direction.
    \item \textbf{Fitness Feedback} ($\mathcal{F}$): A fitness function $\mathcal{F}(P) = \alpha \mathcal{R}_{\text{task}}(P) - \beta \mathcal{C}_{\text{token}}(P) + \gamma \mathcal{I}_{\text{density}}(P)$ that jointly optimizes for task progress, token efficiency, and information density, providing adaptive threshold tuning.
\end{enumerate}

We implement \EpiContext{} as a custom agent within the Harbor evaluation framework~\citep{harbor2026}, a standardized platform for evaluating AI agents in sandboxed environments. This implementation enables direct, fair comparison between \EpiContext{} and baseline context strategies (Full-Context, Sliding Window) under identical conditions, with full result reproducibility.

Our contributions are:
\begin{itemize}
    \item We introduce the epigenetic context model, a framework that applies biological epigenetics concepts to agent context management, establishing a new paradigm where context is actively regulated rather than passively accumulated.
    \item We design three epigenetic operators---methylation, acetylation, and fitness feedback---that together enable dynamic, fitness-driven context evolution without model retraining.
    \item We implement \EpiContext{} as a modular, pluggable agent within the Harbor evaluation framework, demonstrating real-world integration feasibility.
    \item Through controlled experiments comparing six context strategies, we establish baseline performance characteristics and identify the task complexity threshold below which context strategy is irrelevant---a calibration result that defines the applicability boundary for epigenetic context regulation.
    \item We provide a reproducible experimental framework, including agent code, job configurations, raw results, and analysis scripts, for future investigation of context management strategies.
\end{itemize}